\noindent\textbf{Elección del procedimiento.} Con el mismo aire y temperatura, igualar Reynolds exige $\rho_mV_mL_m=\rho_pV_pL_p$. A presión atmosférica las densidades son iguales. Si se limita la velocidad, se eleva la densidad con la presión mediante el gas ideal a temperatura fija.

\begin{enumerate}

\item Calcular el Reynolds que se pretende reproducir.

\[\fitmath{\mathrm{Rep}=\frac{\rho\cdot \mathrm{Vp}\cdot \mathrm{Lp}}{\mu}}\]

\[\fitmath{\mathrm{Rep}=\frac{(1{,}225)\cdot (70)\cdot (2{,}5)}{(1{,}7894\times10^{-5})}\approx 1{,}198027\times10^{7}\,}\]

\item Cancelar densidad y viscosidad en ambos lados y despejar la velocidad del túnel atmosférico.

\[\fitmath{\mathrm{Vatm}=\frac{\mathrm{Vp}\cdot \mathrm{Lp}}{\mathrm{Lm}}}\]

\[\fitmath{\mathrm{Vatm}=\frac{(70)\cdot (2{,}5)}{(0{,}25)}\approx 700\,\mathrm{m/s}}\]

\item El Mach del ensayo limitado usa 80 m/s.

\[\fitmath{\mathrm{Mam}=\frac{\mathrm{Vm}}{\mathrm{asound}}}\]

\[\fitmath{\mathrm{Mam}=\frac{(80)}{(340)}\approx 0{,}2352941\,}\]

\item A temperatura fija $\rho_m/\rho_p=p_m/p_0$; despejar la presión absoluta de ensayo.

\[\fitmath{\mathrm{pm}=\frac{p_{0}\cdot \mathrm{Vp}\cdot \mathrm{Lp}}{\mathrm{Vm}\cdot \mathrm{Lm}}}\]

\[\fitmath{\mathrm{pm}=\frac{(101325)\cdot (70)\cdot (2{,}5)}{(80)\cdot (0{,}25)}\approx 886593{,}8\,\mathrm{Pa}}\]

\end{enumerate}
